\subsection{导出列与下中心列}

\([\mathfrak{g},\mathfrak{g}]\) 这一特征理想称为李代数 \(\mathfrak{g}\) 的导出理想，记作 \(\mathcal{D}\mathfrak{g}\)。

包含 \(\mathcal{D}\mathfrak{g}\) 的 g 的每个子模都是 g 的理想。

\(\mathfrak{g}\) 的导出列是递减序列 \(\mathcal{D}^0\mathfrak{g},\mathcal{D}^1\mathfrak{g},\ldots\)，其各项是 \(\mathfrak{g}\) 的特征理想，递归定义如下：(1) \(\mathcal{D}^0\mathfrak{g} = \mathfrak{g}\)；(2) \(\mathcal{D}^{p + 1}\mathfrak{g} = [\mathcal{D}^p\mathfrak{g},\mathcal{D}^p\mathfrak{g}]\)。

g 的下中心列是 g 的特征理想构成的递减序列 \(\mathcal{C}^1\mathfrak{g},\mathcal{C}^2\mathfrak{g},\ldots\)，递归定义如下：(1) \(\mathcal{C}^1\mathfrak{g} = \mathfrak{g}\)；(2) \(\mathcal{C}^{p + 1}\mathfrak{g} = [\mathfrak{g},\mathcal{C}^p\mathfrak{g}]\)。于是 \(\mathcal{C}^2\mathfrak{g} = \mathcal{D}\mathfrak{g}\)，并且 \(\mathcal{C}^{p + 1}\mathfrak{g}\supset \mathcal{D}^{p}\mathfrak{g}\) 对所有 \(p\) 都成立，这由对 \(p\) 的归纳立即可见。

命题 4. 设 \(\mathfrak{g}\) 和 \(\mathfrak{h}\) 是 \(\mathbf{K}\) 上的两个李代数，且 \(f\) 是从 \(\mathfrak{g}\) 到 \(\mathfrak{h}\) 上的满同态。则 \(f(\mathcal{D}^p\mathfrak{g}) = \mathcal{D}^{p}\mathfrak{h}, f(\mathcal{C}^p\mathfrak{g}) = \mathcal{C}^{p}\mathfrak{h}\)。

若 a 和 b 是 g 的子模，则立即得到

\[
f ([ \mathfrak {a}, \mathfrak {b} ]) = [ f (\mathfrak {a}), f (\mathfrak {b}) ].
\]

于是命题由对 p 的归纳立即得到。

推论. 设 g 是李代数，a 是 g 的理想。李代数 g/a 为交换的充要条件是 a ⊃ \(\mathcal{D}\mathfrak{g}\)。

说 \(\mathfrak{g} / \mathfrak{a}\) 是交换的，就是说 \(\mathcal{D}(\mathfrak{g} / \mathfrak{a}) = \{0\}\)。但根据命题 4，\(\mathcal{D}(\mathfrak{g} / \mathfrak{a})\) 正是 \(\mathcal{D}\mathfrak{g}\) 在 \(\mathfrak{g} / \mathfrak{a}\) 中的典范像。

